% Introduction of main.tex as of 2026-10-14, preserved before the rewrite. Not compiled.
\section{Introduction}\label{sec:introduction}

Vision-language-action (VLA) driving models increasingly generate natural-language reasoning alongside driving plans~\cite{r1,r2,r3}. Reasoning, however, does not guarantee acceptable behavior. An analysis of 150 reasoning--trajectory pairs from Alpamayo 1.5 judged 33.3\% of the reasoning traces unreliable and found 74\% trajectory agreement even for reliable ones~\cite{r18}. A further gap remains when reasoning and trajectory do agree: behavior must also satisfy externally specified constraints, such as traffic rules or operator instructions, that the model's own explanation need not encode. Chain of Causation (CoC) in Alpamayo-R1 illustrates this gap~\cite{r1}. CoC explains why observations lead to a decision, as in ``a pedestrian is crossing, so the vehicle slows down.'' A constraint such as ``do not enter while the zone is occupied, and wait 1.5 s after it clears'' instead restricts which actions are admissible. Because the same scene admits different entry times under different waiting rules, a correct explanation cannot establish compliance.\par

Simply stating the rule to the model is not sufficient. When we supplied explicit speed limits to Alpamayo-R1-10B at inference, choosing limits that the current speed already satisfied, predictions still exceeded them within two seconds in 6/9 PhysicalAI~\cite{nvidia2025physicalai} and 9/9 converted nuScenes~\cite{caesar2020nuscenes} outputs, and raising the limit produced mean XY trajectory differences of at most 0.087 m across inputs and rule-input routes. When asked to hold position within 0.2 m for two seconds, all three predictions for a low-speed input from each source left that radius. The violation counts were the same whether the rule was given in the user prompt or appended to the CoC.\footnote{Alpamayo-R1, also called Alpamayo 1, and the updated Alpamayo 1.5~\cite{r20} belong to the same model family; these experiments use R1 rather than the 1.5 version evaluated in \emph{Drive the Thoughts}~\cite{r18}. Each source contributes three inputs evaluated with three trajectory-generation seeds. The 0.087 m value is the largest three-seed mean XY difference across inputs and rule-input routes over the full 6.4 s prediction horizon; individual paired predictions may differ more. The nuScenes inputs are a third-party conversion to PhysicalAI format; checks use tolerances of 0.05 m/s and 0.01 m. These are open-loop predictions, not measurements of executed driving or road failure rates; the protocol is given in Supplementary Section S1.} Because these diagnostics do not test training, we examined it separately in a controlled synthetic candidate-selection task with a smaller vision-language model (VLM), Qwen3-VL-2B-Instruct~\cite{qwen2025qwen3vl2b}. Starting from the same initialization and giving every adapter the same rule sentence at evaluation, adapters trained on CoC alone violated the rule in 17.92\% of pedestrian-entry predictions, whereas adapters also trained with the constraint violated it in 1.11\%.\footnote{Means over three training seeds on 240 test rows (720 predictions). This is selection among supplied candidates, separate from Alpamayo's trajectory-generation test. Section~\ref{sec:results-training} reports position-hold results and further controls.} An untrained 8B model of the same family, given the same rule, violated it in 68.75\% of these predictions. These results suggest that constraints should enter the training data. This raises a further question: are the constraints themselves correct?\par

Answering this question is difficult for three reasons. First, a single constraint combines evidence with different roles: an observed clearance time, a supplied waiting duration, and an assumption that the zone stays clear. When such a constraint is wrong, a natural-language sentence alone does not reveal which source is at fault. Second, errors are silent but consequential. Constraints that read correctly but were attached to the wrong training cases left accuracy near 50\%, the constraint-free level (Section~\ref{sec:results-association}), and small changes in units, negation, boundaries, or release conditions alter admissible behavior. Third, checking requires a formal representation, whereas a VLM learns from natural language. A checked meaning reaches the model through a sentence, so the sentence must be checked against the specification, and the model's interpretation must still be evaluated: rewording the constraints without changing their meaning reduced static-task accuracy of the CNL-trained model from 100\% to 80.32\% (Section~\ref{sec:results-robustness}).\par

Existing approaches offer complementary foundations. Traffic-rule formalization and verification, regulation retrieval, and safety filtering address rule semantics and use~\cite{r6,Wan2024,r7,r19}. Natural-language formalization connects requirements with temporal logic~\cite{r8,r9}. Natural-language principle learning uses policies to guide model training~\cite{Bai2022,Guan2024}. Our earlier natural-language CoC guards improved action selection but left their construction unchecked~\cite{r4}.\par

We present GuardSynth, a method for developing checkable behavioral constraints for training data, with one component for each difficulty. (i) To separate evidence roles, it binds constraint fields to their sources: observations, supplied rules, and modeling assumptions. (ii) To expose silent errors, it records applicability, action restrictions, persistence and release, timing, and evidence in Evidence-Bound Lifecycle Contracts (EBLC). Each supported contract profile is lowered to Satisfiability Modulo Theories (SMT) queries that check consistency and designated properties within a restricted logical model. (iii) To connect the checked meaning to learning inputs, it generates Controlled Natural Language (CNL) from the checked fields, uses supported fixed-grammar parsers to check that the sentences recover those fields, and appends CNL to the original CoC; model interpretation is evaluated separately. For real scenes, the required observations and applicability need human confirmation.\par

We evaluate GuardSynth with Qwen3-VL-2B-Instruct and low-rank adaptation (LoRA)~\cite{r14} on synthetic action-selection tasks over three seeds. In a fresh-adapter study with identical initialization and evaluation input, training with CNL under cross-entropy loss alone reduced violations from 17.92\%, 6.51\%, and 6.25\% to 1.11\%, 0.00\%, and 1.56\% for pedestrian entry and two position-hold representations, and raised rule-pair accuracy from at most 56.77\% to at least 86.98\%. A semantic loss without rule text reduced violations but left accuracy at the rule-free ceiling, and an untrained 8B model given the rule violated it frequently. Across six tasks, natural-language constraints (P1) and EBLC-derived CNL (P2) let models change their choices with the rule, whereas shuffled case--constraint associations left accuracy near 50\%. The behavioral gain stems mainly from supplying correctly associated constraints, and P2 does not outperform P1 under every condition. What GuardSynth adds is checkability. A fault-injection evaluation targets the pedestrian-entry and position-hold specifications used for learning. Reusing their validators, logical lowering, and renderers with added source- and sentence-correspondence checks, the evaluation detects all 288 injected errors of six supported types and accepts all 48 valid cases. It misses 48 jointly wrong evidence--specification cases. These additional correspondence checks were introduced for this evaluation, rather than applied during the reported training runs (Section~\ref{sec:method-checks}).\par

Our contributions are:\par

\begin{enumerate}
\item A procedure that selects supplied behavioral constraints relevant to a CoC example and binds each constraint field to an observation, a supplied rule, or a modeling assumption.
\item An EBLC-based path from evidence-bound specifications through restricted SMT checking to CNL appended to CoC.
\item An evaluation that separates the model's use of supplied constraints from the quality of the constraints themselves. Under identical initialization and evaluation input, training with constraint text reduces violations and improves rule-dependent selection, which inference-time rules, a larger untrained model, and a rule-free semantic loss do not reproduce. Correctly associated constraints account for most of the behavioral gain, learned behavior is sensitive to wording, and fault injection into the two learning specifications detects supported errors but not jointly wrong evidence and specification.
\end{enumerate}

Section~\ref{sec:related-work} reviews related work, Section~\ref{sec:method} presents GuardSynth, and Sections~\ref{sec:experimental-setup} and~\ref{sec:results} describe the experimental setup and results. Section~\ref{sec:discussion} discusses implications and limitations, and Section~\ref{sec:conclusion} concludes.\par
